% tu-delft-poster — an A0 portrait poster in TU Delft's corporate
% colours, on tikzposter (ships with TeX Live). See README.md for the
% font substitution and class-substitution notes.
\documentclass[a0paper, portrait]{tikzposter}

\usepackage[utf8]{inputenc}
\usepackage[T1]{fontenc}
\usepackage{graphicx}
\usepackage{booktabs}
\usepackage{lipsum}

% TU Delft's own corporate design specifies Roboto Slab (headers) and
% Arial (body). Neither resolves by fontconfig family name in the
% pinned compiler image — only DejaVu and URW base-35 families are
% reachable there — so \setmainfont with either name would silently
% render blank text. Substituted with the closest available families:
% DejaVu Serif in place of Roboto Slab, DejaVu Sans in place of Arial.
\renewcommand{\familydefault}{\sfdefault}

% TU Delft corporate palette (RGB, from tudelft.nl/en/tu-delft-corporate-design).
\definecolor{tudelftblue}{RGB}{0,166,214}
\definecolor{tudelftdarkblue}{RGB}{0,31,77}
\definecolor{tudelftblack}{RGB}{0,0,0}

\definecolorstyle{tudelft}{
  \colorlet{colorOne}{tudelftdarkblue}
  \colorlet{colorTwo}{tudelftblue}
  \colorlet{colorThree}{white}
}{
  \colorlet{backgroundcolor}{white}
  \colorlet{framecolor}{tudelftdarkblue}
  \colorlet{titlefgcolor}{white}
  \colorlet{titlebgcolor}{tudelftdarkblue}
  \colorlet{blocktitlebgcolor}{tudelftblue}
  \colorlet{blocktitlefgcolor}{white}
  \colorlet{blockbodybgcolor}{white}
  \colorlet{blockbodyfgcolor}{black}
}
\usecolorstyle{tudelft}

\title{Poster Title Goes Here}
\author{Author One, Author Two \quad\textbar\quad Faculty Name, TU Delft}
\institute{\includegraphics[height=2.5cm]{tudelft-logo}}

\begin{document}

\maketitle

\begin{columns}
\column{0.33}

\block{Introduction}{
\lipsum[1]

\lipsum[2]

\lipsum[10]
}

\block{Related Work}{
\lipsum[3]

\lipsum[11]

\begin{itemize}
\item A first strand of prior work addressing a related but narrower
question.
\item A second strand that shares this poster's method but applies it
to a different domain.
\item A gap the present work closes, stated as concretely as possible.
\end{itemize}
}

\block{Research Questions}{
\begin{itemize}
\item How does the proposed method compare against the strongest
existing baseline?
\item Does the improvement hold across different data regimes?
\item What is the computational cost of the improvement?
\end{itemize}
}

\column{0.34}

\block{Method}{
\lipsum[4]

A results table:
\vspace{0.5em}

\begin{tabular}{lr}
\hline
Metric & Value \\
\hline
Accuracy & 0.94 \\
Precision & 0.91 \\
Recall & 0.89 \\
Latency (ms) & 12 \\
\hline
\end{tabular}

\vspace{1em}
\lipsum[5]

\lipsum[12]
}

\block{Implementation Notes}{
\lipsum[6]

\lipsum[13]
}

\block{Limitations}{
\lipsum[14]
}

\column{0.33}

\block{Results}{
\lipsum[7]

\lipsum[8]

\lipsum[15]
}

\block{Discussion}{
\lipsum[16]
}

\block{Conclusion}{
\lipsum[9]

Future work will extend this approach along the following directions:
\begin{itemize}
\item Scaling the method to larger, noisier datasets.
\item Evaluating generalisation across institutions.
\item Releasing the implementation as an open-source tool.
\end{itemize}
}

\block{Acknowledgements}{
\lipsum[17]
}

\end{columns}

\block{Contact}{
\texttt{author@tudelft.nl} \hfill TU Delft corporate design used with
permission of the project owner.
}

\end{document}
